\section*{Chapter \ref{ch:b2:adverse-selection} --- Adverse Selection}

\begin{solution}{exo:b2:adverse-selection:1}
% code: exercises/solutions.py lemons
The sellers whose goods are worth at most 6{,}000, 60 per cent of them. The average quality on offer is 3{,}000, worth $1.2 \times 3{,}000 =
3{,}600$ to a buyer, less than the price: buyers do not pay 6{,}000.
\end{solution}

\begin{solution}{exo:b2:adverse-selection:2}
% figdata: annuities_summary
% ledger: b2:fp2004-t2
Buyers of index-linked and escalating annuities lived longer; buyers of guaranteed and capital-protected annuities died
sooner. The effects are larger in the voluntary market: 18.5 and 15.2 points fewer deaths within five years, against 5.3
and 7.2 in the compulsory market.
\end{solution}

\begin{solution}{exo:b2:adverse-selection:3}
% code: exercises/solutions.py pooling_threshold
A pooled premium with a share $s$ of low risks is $1{,}000s + 3{,}000(1 - s)$. Low risks buy if it is at most 1{,}500: $3{,}000 - 2{,}000s \le 1{,}500$, so $s \ge 0.75$.
\end{solution}

\begin{solution}{exo:b2:adverse-selection:4}
Moral hazard: drivers with more cover may take less care. Characteristics that the insurer prices on: drivers who drive more,
or in cities, may both buy more cover and have more accidents. Comparing drivers whom the insurer prices alike removes the
second; separating the first needs variation in cover that drivers do not choose, or the timing of claims after a change
in cover.
\end{solution}

\begin{solution}{exo:b2:adverse-selection:5}
At a price $p \le 1$ the average quality offered is $p/2$, worth $bp/2$; buyers pay $p$ if $bp/2 \ge p$, that is, if $b \ge 2$. With $b \ge 2$ there is trade
at any price up to 1; with $b < 2$ there is none.
\end{solution}

\begin{solution}{exo:b2:adverse-selection:6}
% code: exercises/solutions.py spiral
At 596, $(900 - 596)/500 = 0.608$ buy; their average cost is $800 - 300 \times 0.608 = 617.6$. At 617.6, 0.565 buy, with average cost 630.6, at which
0.539 buy.
\end{solution}

\begin{solution}{exo:b2:adverse-selection:7}
% code: exercises/solutions.py linear_market
$\mathit{AC}(q) = \frac1q\int_0^q (800 - 600x)\,dx = 800 - 300q$. Equilibrium: $900 - 500q = 800 - 300q$, so $q = 0.5$ and the premium is 650. Since $D - \mathit{MC} = 100 +
100q > 0$, everyone should be insured; the loss is $\int_{0.5}^1 (100 + 100q)\,dq = 50 + 37.5 = 87.5$ dalens a person.
\end{solution}

\begin{solution}{exo:b2:adverse-selection:8}
% code: exercises/solutions.py linear_market
$\mathit{AC}(q) = 300 + 100q$. Demand is at least average cost for every $q$: at $q = 1$, $400 = 400$, so everyone buys at a premium of 400. Efficient
coverage is where $900 - 500q = 300 + 200q$, so $q = 6/7 = 0.857$. The members beyond it value cover below its cost; the loss is $\int_{6/7}^1 (700q
- 600)\,dq = 7.14$ dalens a person: too much insurance.
\end{solution}

\begin{solution}{exo:b2:adverse-selection:9}
Each product was priced for the average buyer of a given age and sex, but buyers knew more about their own health and chose the
product that paid best for it. Pricing each feature on its own buyers' mortality would end the subsidy from one product to
another, but within each product buyers would still differ in what they knew, so selection on whether to buy at all would
remain.
\end{solution}

\begin{solution}{exo:b2:adverse-selection:10}
At $p \le 1$ buyers value the goods offered at $bp/2 < p$ since $b < 2$; above 1, at $b/2 < 1 < p$. So no positive price clears the market. With full
information every good would be traded, with gains $\int_0^1 (b - 1)\theta\,d\theta = (b - 1)/2$, all lost.
\end{solution}

\begin{solution}{exo:b2:adverse-selection:11}
$\mathit{AC}(q) = m_0 - m_1q/2$. Equilibrium: $d_0 - d_1q = m_0 - m_1q/2$, so $q^* = (d_0 - m_0)/(d_1 - m_1/2)$ when $d_1 > m_1/2$, capped at 1. If $d_0 \le m_0$, demand lies at or
below average cost at every positive share: even the person who values cover most would not pay the average cost of
insuring her alone, and the market unravels completely.
\end{solution}

\begin{solution}{exo:b2:adverse-selection:12}
% code: exercises/solutions.py linear_market
With everyone insured the premium is $\mathit{AC}(1) = 500$. The person who values cover least would pay $D(1) = 400$; she buys if the penalty makes
staying out cost more than the gap, so the penalty must be at least 100 dalens.
\end{solution}

\begin{solution}{pb:b2:adverse-selection:1}
% code: exercises/solutions.py coop
\textbf{Part I.}
1. $(600 \times 400 + 400 \times 1{,}600)/1{,}000 = 880$ kests.
2. $600 \times 200 + 400 \times 300 = 240{,}000$.
3. Yes: each type values cover above its cost.
4. 400 and 1{,}600.
5. 200 for the healthy, 300 for the sick.
6. Each member's health risk, hidden from the cooperative.
\textbf{Part II.}
7. Only the sick: 600 is below 880.
8. 1{,}600; the sick buy.
9. 400.
10. The healthy members' surplus, $600 \times 200 = 120{,}000$.
11. $400 \times 300 = 120{,}000$.
12. No: breaking even with everyone insured needs 880, more than the healthy will pay.
\textbf{Part III.}
13. $1{,}500 - 600q$.
14. $1{,}800 - 1{,}000q = 1{,}500 - 600q$: $q = 0.75$, premium 1{,}050.
15. All members: $1{,}800 - 1{,}000q$ exceeds $1{,}500 - 1{,}200q$ by $300 + 200q$.
16. $1{,}000 \times \int_{0.75}^1 (300 + 200q)\,dq = 1{,}000 \times (75 + 43.75) = 118{,}750$ kests.
17. At 900, 0.9 buy, with average cost 960; at 960, 0.84 buy, with average cost 996; at 996, 0.804 buy.
18. Demand stays above average cost up to 0.75, so the members who remain there will pay the average cost of the group.
\textbf{Part IV.}
19. $\mathit{AC}(1) = 900$; the gain is the 118{,}750 lost in Part III.
20. Members ranked above 0.9, whose willingness to pay is below 900: 100 members.
21. 100 kests: the gap between 900 and the lowest willingness to pay, 800.
22. Everyone, since each values cover above its cost; the highest-cost member pays 1{,}500.
23. 1{,}050 with no intervention, 900 under the mandate, 1{,}500 under risk rating.
24. Without intervention 75 per cent of members are insured at 1{,}050 kests and selection costs 118{,}750 kests a year; a mandate with a
penalty of at least 100 insures everyone at 900; risk-rated premiums insure everyone at premiums from 300 to 1{,}500, the sickest
paying most.
\end{solution}
